\chapter{Korrekturregister und ergänzende Details}
\label{app:corrections}
\begin{longtable}{L{46mm}L{77mm}L{27mm}}
\toprule \textbf{Frühere Formulierung}&\textbf{Konsolidierte Aussage}&\textbf{Herkunft}\\\midrule\endhead
Aus zwei Zahlen folgt der ganze Aufbau.&P1/P2 enthalten Orientierung, Positivität, Trägerart und Markierungen. Ihre Konsequenzen sind von ihrer Auswahl verschieden.&S09; S05/S06\\
Die Kontextregel ist symmetrieerzwungen.&$B/7$ ist eine konkrete gleichgewichtete Realisierung; die kovariante Familie $K_a$ bleibt ein Gegenbeispiel zur Eindeutigkeit.&S05/S06\\
Es existiert kein autonomer Systemschatten.&Im festen Strahlenprozess gilt $FT=(3/7)F$. Nichtautonomie entsteht relativ zum erweiterten Registerzugriff.&S01; S06\\
240 CQ-Koordinaten sind die minimale Hülle.&Für verschiedene passive Klassen gelten 240, 60, 45 oder 30 lineare Dimensionen. Kohärente Register sind im Allgemeinen nicht CQ.&S01; S06\\
30 Nullrichtungen speichern Gedächtnis.&Sie werden durch $T$ ausgelöscht. Das relevante Gedächtnis sitzt in größeren zugänglichen Registern.&S07; S01\\
Die 240-Zählung identifiziert $E_8$.&Wurzeln, Koordinaten und Ankerprodukte sind verschiedene Objekte. Die natürliche äquivariante Identifikation wird in S11 zusätzlich obstruiert.&S01; S11\\
Die 248 verbietet symmetrische Paarzustände.&Ein Generatorsektor ist kein allgemeiner Mehrteilchen-Zustandsraum. Eine harte Hülle ist eine Zusatzbedingung.&S01/S05/S06\\
Genau vier Träger sind allein daraus erzwungen.&Vier ist maximal und singulär unter der harten Antisymmetriebedingung; kleinere antisymmetrische Räume existieren.&S01/S05\\
Ein bekanntes $\Omega$ wählt den Graphen.&Viele verbundene positive Vierergraphen haben denselben Kern und verschiedene Spektren.&S01/S05/S06\\
Der vollständige $4m$-Graph hat immer genau einen Grundzustand.&$E_0=5m(m-1)J$ und Gap $2J$, aber Grundraummultiplizitäten $1,14,462,24024$ für $N=4,8,12,16$.&Hier ergänzt\\
Gemischter Spinorkanal führt direkt nach $(10,6)$.&$[\mathfrak g_1,\mathfrak g_3]\subset\mathfrak g_0$; nur $1+1\to2$ liefert diesen direkten Zielgrad.&S01/S05/\allowbreak S06/S11\\
Fehlende Wurzelsumme erzwingt hard-core.&Sie verbietet den direkten Klammerkanal. Die physische Besetzungsregel benötigt eine Fock-/Zustandszuordnung.&Hier präzisiert\\
Clebsch-Uniformität erzwingt die kritische Kette.&Der Graph hat gewechselt. Uniforme Gewichte auf Clebsch sind keine Herleitung einer eindimensionalen Kette.&Hier präzisiert\\
Aufzeichnung benötigt zwingend einen neuen $E_8$-Modul.&Die konkrete Kopplung ist ein Polynom in trägerübergreifenden lokalen Matrixgeneratoren. Die primitive Zugriffserlaubnis bleibt offen.&Hier ergänzt\\
Die Tetramer-Paarenergien bilden ein freies Uhrwerk.&Sie kommutieren mit $H_{\mathrm{tet}}$. Laufende andere Observablen und bedingte Viereruhren sind möglich.&S01; S03; S05\\
Ein Singulett präpariert sich durch Austausch.&Der Singulettanteil ist unter reinem Austausch erhalten. S02 ergänzt ausdrücklich eine andere Ressource und Erfolgsselektion.&S02/S03\\
Der Rückkehrwert ist die Roh-Erfolgsrate.&Endfilterfaktor $9/16$ und Präparationsfaktor $3/32$ müssen mitgeführt werden.&S02\\
$c=8$ ist ein neuer Kondo-Fixpunktbeweis.&Das Produkt hat bereits $5+3=8$. Die konkrete Aufgabe ist die geeignete lokale Erweiterung und ihre native Realisierung.&S01/S11\\
Eine Ringkennzahl beweist die chirale Naht.&Linke und rechte Sektoren, $D_5$, Glue, Operatorprodukte und native analytische Kontrolle bleiben nötig.&S01/S05/S06\\
Glue-$\Z_4$ ist geometrische Vierteldrehung.&Die Spuren auf Stufe eins sind null und $248i$. Die Operatoren sind verschieden.&S08\\
Unendliche Register leiten den Zeitpfeil ab.&Der eingehende Zustand, Zugriff und die Korrelationsstruktur sind zusätzliche Auswahlbedingungen.&S01/S05\\
Ein genauer Formelwert ist ein abgeschlossener Physiknachweis.&Quellgleichung, Transfer, empirische Unsicherheit und negative Ausgaben müssen gemeinsam geprüft werden.&Alle Audits\\\bottomrule
\end{longtable}

\section{Frühere Nebenstränge, die im Gesamtbild erhalten bleiben}
S09 berichtet für ein früheres Majoranamodell zehn zugängliche und sechs verborgene reelle Richtungen, entsprechend einer $32\times8$-Zerlegung auf Fockebene. Verschränkung allein macht einen verborgenen Zustand nicht am Rand lesbar; eine geeignete Kopplung oder erweiterte Messung ist nötig. Ein dort gebauter Ereigniskanal übertrug Besetzung, besaß aber weiterhin eine nichttriviale Fixalgebra. Übertragung und eindeutige Zustandsauswahl sind verschiedene Leistungen.

Unter früheren lokalen Symmetrieannahmen blieben sieben Kanalparameter: vier lokal bestimmbar, drei erst durch gemeinsame Antworten. Eine stärkere Symmetrie führte später zum Kanaldreieck. Eine engere Vier-Generatoren-Ereignishypothese lieferte bei gleichen Raten das binomiale Muster $1$--$4$--$6$--$4$--$1$. Diese verschiedenen Parameterzahlen widersprechen sich nicht; sie verwenden unterschiedliche Eingaben. Keine davon wählt ihre Ereignisse und Raten von selbst aus.

Der Ikosaeder-/McKay-Vergleich aus S09 ist ein weiterer Strukturhinweis: Eine gewöhnliche $A_5$-Darstellung bleibt beim Hochziehen zur binären Ikosaedergruppe zentral trivial. Der zweidimensionale Spinorsektor, in dem das zentrale Element als $-I$ wirkt, entsteht dadurch nicht automatisch. Eine ähnliche Ikosaedersymmetrie einer anderen Theorie schließt daher keinen fehlenden spinoriellen Quellsektor.

Die in S04 genannte Taxonomie gescheiterter Forschungsrouten enthält Zirkularität, verlorene Information, fehlende Brücken, Strukturfehler, numerische Artefakte, implizite Orakel, bloße Umformulierungen, fehlende Weltabhängigkeit und unkontrollierte Grenzwerte. Diese Kategorien sind nützlich, weil sie eine gescheiterte Schlussfolgerung präzise benennen. Die zugehörigen alten Zählungen sind kein aktueller Audit dieses Buches.

\chapter{Begriffe in einfachen Worten}
\begin{longtable}{L{34mm}L{116mm}}
\toprule \textbf{Begriff}&\textbf{Bedeutung in diesem Dokument}\\\midrule\endhead
Adjungierte&Die Darstellung einer Lie-Algebra auf ihren eigenen Generatoren durch die Klammer; nicht automatisch der gesamte Raum physischer Zustände.\\
Adjunktion $\dagger$&Komplex konjugierte Transposition beziehungsweise die passende Operatoradjunktion. Nicht dasselbe Wort wie adjungierte Darstellung.\\
Algebra&Ein Raum von Objekten mit Regeln für Addition, Multiplikation und gegebenenfalls Adjunktion.\\
Antisymmetrie&Vertauschen zweier Plätze ändert das Vorzeichen.\\
Anregung&Ein Zustand oberhalb des Grundzustands, der Energie und unter geeigneten Bedingungen Information oder Ladung tragen kann.\\
Chiralität&Unterschied zwischen links- und rechtshändigen beziehungsweise links- und rechtslaufenden Sektoren; die genaue Bedeutung hängt von der Dimension ab.\\
Clifford/Stabilizer&Ein spezieller Quantenbaukasten mit einfacher Pauli-Struktur. Unter geeigneten Präparations- und Messbedingungen klassisch effizient simulierbar.\\
Compiler&Hier eine festgelegte Übersetzung zwischen mathematischen Beschreibungen. Kein bereits gefundenes kosmisches Computerprogramm.\\
CQ-Zustand&Klassisches Kontextlabel mit zugehörigen Quantenblöcken. Zwischen verschiedenen klassischen Labels gibt es keine kohärenten Offdiagonalen.\\
Dephasierung&Verlust bestimmter relativer Phasen in einer reduzierten Beschreibung; diagonal betrachtete Wahrscheinlichkeiten bleiben erhalten.\\
Energie-Gap&Abstand zwischen Grundenergie und nächster Anregung. Ein endlicher Gap ist noch kein thermodynamischer Beweis.\\
Fockraum&Ein Zustandsraum mit verschiedenen Besetzungszahlen von Moden und ausdrücklich festgelegter Teilchenstatistik.\\
Glue / Verklebung&Kompatible Ergänzung verschiedener Ladungs- oder Gitterrestklassen zu einem größeren gemeinsamen Objekt.\\
Hamiltonoperator&Der Operator, der Energien und bei gegebener Zeitstruktur die unitäre Entwicklung bestimmt.\\
Heraldierte Präparation&Ein Herstellungsversuch mit eigener Anzeige, ob der gewünschte Zustand erfolgreich erzeugt wurde.\\
Inzidenz&Die genaue Angabe, welche Kontexte oder geometrischen Objekte gemeinsame Bestandteile besitzen.\\
Kohärenz&Erhaltene relative Phasen, die spätere Interferenz ermöglichen.\\
Kontext&Eine vollständige Menge gemeinsam messbarer Alternativen; hier eine orthogonale Viererbasis.\\
Krausoperator&Eine mathematische Beschreibung eines möglichen Zweigs einer Quantenoperation.\\
Lift&Eine reichhaltigere Ausführung, die unter einer eingeschränkten Auslesung dieselbe reduzierte Regel ergibt.\\
Marginale&Der Zustand oder die Statistik eines Teils, wenn man den Rest nicht mitausliest.\\
Modularer Fluss&Eine aus Algebra und Zustand gebaute innere Entwicklung; in endlichen Beispielen durch Potenzen von $\rho$ beschrieben. Nicht automatisch physikalische Zeit.\\
Operationelle Gleichheit&Kein erlaubtes Zukunftsexperiment kann zwei Beschreibungen unterscheiden.\\
Pointer / Register&Ein Hilfssystem, in dem eine Aufzeichnung gespeichert wird. Es kann kohärent, dephasiert, frisch oder wiederverwendet sein.\\
Positive Statistik&Wahrscheinlichkeiten und Korrelationsmatrizen erfüllen die nötigen Positivitäts- und Normierungsbedingungen.\\
Propagator&Eine Korrelationsfunktion, die beschreibt, wie eine Anregung zwischen Ereignissen oder Orten übertragen wird.\\
Ququart&Ein Quantensystem mit vier Basiszuständen; hier meist durch zwei Qubits dargestellt.\\
Schatten&Eine mathematisch festgelegte Auslesung, die einen Teil der ursprünglichen Information verwirft.\\
Singulett&Ein Zustand, auf den die betreffende Symmetrie trivial wirkt. Ein großer Raum kann viele Kopien eines Singuletts enthalten.\\
Strahl&Ein reiner Quantenzustand ohne seine globale Phase.\\
Swap&Vertauschung zweier Träger. Der Operator hat symmetrische und antisymmetrische Eigenräume.\\
WZW / zentrale Ladung&Eine Klasse konformer Feldtheorien und eine ihrer Kennzahlen. Eine passende Kennzahl allein bestimmt nicht alle Operatoren oder die Händigkeit.\\\bottomrule
\end{longtable}

\chapter{Quellen, Abdeckung und Reproduktion}
\section{Die elf Eingabedokumente}
Die Quellenkürzel S01 bis S11 gelten durchgängig für dieses Buch und sind von den jeweils eigenen Kürzeln der Anhänge zu unterscheiden. Die vollständigen Originalpfade und SHA-256-Werte stehen im beiliegenden Manifest. Hier werden die Dateinamen und inhaltlichen Rollen aufgeführt.
\input{sources_manifest}

\section{Was neu geprüft wurde}
Der eigenständige Prüfer importiert keinen Code des ursprünglichen TFPT-Repositoriums. Die Standard-Zwei-Qubit-Kontexte wurden aus den üblichen Pauli-Matrizen rekonstruiert. Das kontrolliert die endlichen im Text angegebenen Beziehungen, ist aber kein erneuter phasentreuer Import des nativen Gauß-$E_8$-Adapters.

Die 210 Bedingungen umfassen insbesondere:
\begin{itemize}
\item 15 Kontexte, 60 Projektoren, Pauli-Gramidentitäten, die vollständige $T$-Faktorisierung gegen Born-Überlappungen, Ränge 30/45/60 und die CQ-Fortsetzung;
\item den Clebsch-Graphen aus den 16 Spinorgewichten, stark reguläre Identität, charakteristisches Polynom, zehn Bindungslabels und fünf perfekte Matchings;
\item alle Schur-Weyl-Sektoren der vollständigen SU(4)-Graphen bei $N=4,8,12,16$, inklusive Grundraummultiplizität und Gesamt-Hilbertraumdimension;
\item den ganzzahligen $\Omega$-Tensor, Swap-Pauli-Identität, Paarmarginale, Aufzeichnungsprojektor und seine Wirkung;
\item das originale und das ausgeglichene Echo direkt auf den Tensoramplituden, einschließlich positiver und negativer Kontrolle durch den Tick beziehungsweise seinen Wegfall;
\item den invarianten Zweizellenraum auf allen 65536 Komponenten, symbolische Eigenwert- und Reihenidentitäten;
\item die Hermitesche-$2\times2$-Determinante und eine hochpräzise numerische Nullstelle der elektromagnetischen Gleichung.
\end{itemize}
Die kleinen Pauli- und Projektorrechnungen verwenden ganze oder dyadische Werte, die in den verwendeten Operationen exakt darstellbar sind; Ränge, Hookformeln und Polynomidentitäten werden rational oder symbolisch geprüft. Die Dezimalwurzel bleibt eine numerische Kontrolle. Die Diagrammkurven stellen die angegebenen Formeln dar und sind keine zusätzlichen Messungen.

Die Original-Clebsch-Vielteilchendiagonalisierung, die vollständigen Ringrechnungen, Zweischleifenläufe, 3875-Verzweigung, ursprünglichen Quelladapter-Prüfungen und die 20-Qubit-Schaltungen wurden für dieses Buch nicht neu ausgeführt. Ihre Herkunft wird bei den entsprechenden Aussagen genannt. Ein behaupteter grüner Originallauf wird nicht durch das bloße Lesen seiner Zusammenfassung zu einem hier neu geprüften Lauf.

\section{Wie die Quellen ihre eigenen Prüfungen beschreiben}
\par\noindent\begin{tabularx}{\linewidth}{L{16mm}L{41mm}Y}
\toprule \textbf{Quelle}&\textbf{Berichtete Prüfung}&\textbf{Reichweite}\\\midrule
S09&11214 Bedingungen pro Modus, elf Negativvarianten&Ursprungsaudits und Zahlenprüfungen; kein vollständiger Lean-/RH-Neubau.\\
S07&325 + 147 Bedingungen, 19 Tests&Inversion und Follow-ups; mehrere spätere Minimalitätskorrekturen erforderlich.\\
S06&86 eigene Bedingungen sowie ausgewählte Originalreplays&45-/30-Quotienten, Uhr und Herkunftskontrollen; begrenzte Kettenklasse.\\
S01&148 Bedingungen, davon 108 exakt und 40 numerisch&Unabhängiger endlicher Nachbau ohne Originalrepository.\\
S02&598 benannte Bedingungen&Schaltungs- und Rohwahrscheinlichkeitsprüfungen; keine Hardware.\\
S11&27288 Bedingungen, 14 Tests; zusätzliche Schur-Weyl-Numerik&Graph-/Sektorprüfungen und numerische Clusterwerte; Herkunftslesart bleibt bedingt.\\\bottomrule
\end{tabularx}\par
Die Zahlen werden nicht addiert: Die Prüfungen teilen Quellobjekte und wiederholen Identitäten. Sie sind keine Zahl unabhängiger physikalischer Bestätigungen.

\section{Reproduktionspaket dieser Synthese}
Das Paket enthält die vollständige LaTeX-Quelle, Kapiteldateien, Vektordiagramme, den eigenständigen Prüfer, die Ergebnisdatei und das Eingabemanifest. Mit einer LaTeX-Installation einschließlich LuaLaTeX, TikZ und den verwendeten Standardpaketen lässt sich das Buch durch zwei Läufe von \texttt{lualatex main.tex} im Quellverzeichnis bauen. Für die endlichen Gegenrechnungen werden Python, NumPy, SymPy, mpmath und Matplotlib benötigt.

\texttt{python3 verify\_synthesis.py} erzeugt die Ergebnisdatei und die aus Formeln gebauten Abbildungen. Die Akzeptanzbedingungen verwenden explizite Fehler statt wegoptimierbarer Assertions. Der normale und der mit deaktivierten Assertions gestartete Lauf wurden beide ausgeführt; ihre Ergebnisdateien sind bytegleich. Das Paket enthält nicht sämtliche Originalprüfer der elf Quellen und behauptet daher keine vollständige Reproduktion des ursprünglichen Repository-Stands.

\section{Externe Primärliteratur und Referenzdaten}
Die folgenden Quellen wurden zur fachlichen Einordnung beziehungsweise für den datierten Konstantenvergleich aufgerufen. Sie bestätigen die allgemeinen Werkzeuge, nicht die physikalische Richtigkeit von TFPT. Die spezifischen langen Herleitungen dieses Buches beruhen auf den Anhängen und den ausgeschriebenen Gegenrechnungen.
\begin{enumerate}[label={E\arabic*.},leftmargin=2.4em,itemsep=7pt]
\item F. A. Pollock et al.: \emph{Non-Markovian quantum processes: complete framework and efficient characterisation}. Physical Review A 97, 012127 (2018). \url{https://arxiv.org/abs/1512.00589}
\item S. Bravyi, D. DiVincenzo, D. Loss: \emph{Schrieffer-Wolff transformation for quantum many-body systems}. Annals of Physics 326, 2793--2826 (2011). \url{https://arxiv.org/abs/1105.0675}
\item E. Moreva et al.: \emph{Time from quantum entanglement: an experimental illustration}. \url{https://arxiv.org/abs/1310.4691}
\item S. Bravyi, A. Kitaev: \emph{Universal Quantum Computation with ideal Clifford gates and noisy ancillas}. \url{https://arxiv.org/abs/quant-ph/0403025}
\item M. Führinger et al.: \emph{DMRG studies of critical SU(N) spin chains}. Annalen der Physik 17, 922 (2008). \url{https://arxiv.org/abs/0806.2563}
\item V. G. Kac et al.: \emph{Conformal embeddings and simple current extensions}. IMRN 2015, 5229--5288. \url{https://arxiv.org/abs/1210.6602}
\item G. Chiribella, G. M. D'Ariano, P. Perinotti: \emph{Informational derivation of Quantum Theory}. Physical Review A 84, 012311 (2011). \url{https://arxiv.org/abs/1011.6451}
\item G. Chiribella, G. M. D'Ariano, P. Perinotti: \emph{Theoretical framework for quantum networks}. \url{https://arxiv.org/abs/0904.4483}
\item S. Milz et al.: \emph{Kolmogorov extension theorem for (quantum) causal modelling and general probabilistic theories}. \url{https://arxiv.org/abs/1712.02589}
\item P. Zanardi, D. A. Lidar, S. Lloyd: \emph{Quantum tensor product structures are observable-induced}. Physical Review Letters 92, 060402 (2004). \url{https://arxiv.org/abs/quant-ph/0308043}
\item P. Hořava: \emph{Stability of Fermi Surfaces and K-Theory}. Physical Review Letters 95, 016405 (2005). \url{https://arxiv.org/abs/hep-th/0503006}
\item G. M. D'Ariano, M. Erba, P. Perinotti: \emph{Isotropic quantum walks on lattices and the Weyl equation}. Physical Review A 96, 062101 (2017). \url{https://arxiv.org/abs/1708.00826}
\item A. Connes: \emph{On the spectral characterization of manifolds}. \url{https://arxiv.org/abs/0810.2088}
\item NIST: \emph{2022 CODATA adjustment, Fundamental Physical Constants}. Datierte Referenz, für dieses Buch am 14. September 2026 geprüft. \url{https://physics.nist.gov/cuu/Constants/Table/allascii.txt}
\item I. Esteban et al.: \emph{NuFit-6.0: Updated global analysis of three-flavor neutrino oscillations}. \url{https://arxiv.org/abs/2410.05380}
\item Planck Collaboration: \emph{Planck 2018 results. VI. Cosmological parameters}. \url{https://arxiv.org/abs/1807.06209}
\item K. Zuse: \emph{Rechnender Raum}. Vieweg, 1969. \url{https://link.springer.com/book/10.1007/978-3-663-02723-2}
\item Wolfram Physics Project: \emph{Technical Introduction, Basic Concepts}. Selbstdarstellung des Programms. \url{https://www.wolframphysics.org/technical-introduction/potential-relation-to-physics/basic-concepts/}
\end{enumerate}
\section{Abschluss der Dokumentation}
Die Anhänge wurden unverändert gelesen; ihre tatsächlichen Dateiinhalte sind durch SHA-256 dokumentiert. Dieses Buch aktualisiert keinen ursprünglichen Theorievertrag und keinen Abschlussmarker. Es liefert eine neue, zusammenhängende und prüfbare Darstellung mit ausdrücklich abgegrenzten eigenen Fortsetzungen.
